\section{FOSP melanoma application}\label{sec:fosp-melanoma}

The empirical application considered cutaneous melanoma records from
the Hospital Cancer Registry of S\~ao Paulo (RHC/FOSP), Brazil. The present analysis used
an updated registry extract from July 2026, retaining diagnoses from
2000--2014, clinical stages I--IV, valid positive observation times, and
positive recorded age. The resulting cohort comprised 7,137 patients,
compared with 6,749 in the published application; it is an updated
reconstruction rather than the identical historical analysis sample.
There were 3,626 women (50.8\%), mean age was 57.80 years with standard
deviation 16.33, and 4,864 patients (68.2\%) had stage I or II disease.

The analysis used the extended follow-up fields, with time measured in
years from diagnosis to the last recorded information and death from
cancer as the event of interest. There were 2,239 recorded cancer deaths
and 4,898 censored observations. Deaths from other causes and patients
without a recorded cancer death were treated as censored, giving a
cause-specific analysis rather than an all-cause or competing-risks
analysis. The median and maximum observed times were 7.60 and 25.86 years,
respectively. The alternative fields with administrative censoring in
July 2018, provided for comparison with the earlier application, were not
used in these fits.

The predictor set comprised age at diagnosis, diagnosis year, sex,
clinical stage, primary tumour topography, histological type, and separate
indicators for surgery, radiotherapy, chemotherapy, hormone therapy, and
immunotherapy. Age and diagnosis year were standardized within each
training sample. Clinical stage, topography, and histological type were
represented by smoothed survival-target encodings based on restricted
mean survival time, using five-fold out-of-fold encoding for the training
observations. Validation and test observations received mappings
estimated exclusively from the corresponding training sample. This
followed the training-only preprocessing principle in
Section~\ref{sec:prediction-evaluation}.

Predictive performance was assessed using 50 event-stratified paired
outer holdouts with 70\% training and 30\% testing. All methods shared
the same outer splits, and tuning was confined to three inner training
folds using IPCW integrated Brier score and the horizon rule in
Section~\ref{sec:prediction-evaluation}. The comparison included the seven
GG-KM kernels under Poisson and negative-binomial counts, additive
cubic-spline models under both count laws, and the mixture-cure and
promotion-time comparators. Kernel candidates followed the grids in
Section~\ref{sec:simulation}, including wavelet scales
$A\in\{0.5,1,2,10,15,20\}$. Cubic-spline candidates used three or five
training-fold quantile knots with
$\lambda\in\{0.1,1\}$.

\GGKMTable{sections/generated/fosp-melanoma-performance-table}

Table~\ref{tab:fosp-melanoma-performance} reports the split-level means
and variation. The Poisson and negative-binomial wavelet GG-KM models
attained the lowest mean IPCW integrated Brier score, both 0.1218 with
split standard deviation 0.0036, and the highest mean Uno's C-index,
both 0.8161. The spline models had mean IBS 0.1238, compared with 0.1242
for mixture cure and 0.1248 for promotion time.
The table summarizes 865 successful fits, all of which met the recorded
convergence criteria. The negative-binomial linear model contributed
only 15 successful holdouts, whereas every other specification
contributed 50; its summary therefore represents a smaller subset of
training samples.

Figure~\ref{fig:fosp-melanoma-performance} combines descriptive survival
calibration with the paired prediction-error and survival-discrimination
distributions. The six displayed mean survival curves followed the early
Kaplan--Meier decline and lay slightly below it later in the common
supported horizon. The wavelet models had lower mean IBS than the
classical and spline comparators. Overlapping repeated
holdouts describe variation across training samples and do not provide
independent-sample confidence intervals.

\GGKMWideFigure{fosp_melanoma_predictive_assessment_ibs_uno_c_index.pdf}
{Predictive assessment in 50 paired repeated holdouts of the FOSP melanoma
cohort using extended follow-up. Panel A compares the full-cohort
Kaplan--Meier curve with split-averaged out-of-sample survival predictions
over the common supported horizon. Panel B shows split-level IPCW
integrated Brier scores and Panel C Uno's C-index. Each point represents one
outer split; lower IBS and higher C-index values indicate better
prediction and discrimination, respectively. The selected GG-KM models
use the wavelet kernel.}
{fig:fosp-melanoma-performance}
{Three-panel FOSP melanoma assessment showing the Kaplan--Meier curve and
six mean predicted survival curves, followed by box-and-strip plots of
IBS and Uno's C-index for mixture cure, promotion time, Poisson and
negative-binomial splines, and wavelet GG-KM under both count laws.}

Figure~\ref{fig:fosp-melanoma-performance-ggkm-kernel} shows that kernel
choice materially affected prediction. Wavelet GG-KM had the lowest mean
IBS. The Laplacian models and the negative-binomial polynomial model
followed at 0.1229. The Poisson polynomial and
linear specifications had larger mean errors, 0.1429 and 0.1488.
Poisson and negative-binomial results were close for the RBF, Laplacian,
exponential, Cauchy, and wavelet kernels, whereas the polynomial models
differed substantially. Mean predicted cure probabilities ranged from
0.592 to 0.733 across specifications; these remain model-based
extrapolations rather than observed individual cure status. The results
therefore favor wavelet GG-KM for IBS and concordance in this cohort.

\GGKMWideFigure{fosp_melanoma_ggkm_kernel_comparison_ibs_uno_c_index.pdf}
{Kernel-specific GG-KM performance in the FOSP melanoma cohort. Panel A
shows IPCW integrated Brier scores and Panel B Uno's C-index across seven kernels,
with Poisson and negative-binomial specifications distinguished by color
and hatching. Points represent successful outer holdouts. The
negative-binomial linear model contributes 15 successful fits; all other
kernel/count combinations contribute 50.}
{fig:fosp-melanoma-performance-ggkm-kernel}
{Two-panel comparison of split-level IBS and Uno's C-index for linear, RBF,
Laplacian, exponential, Cauchy, polynomial, and wavelet GG-KM models under
Poisson and negative-binomial counts in the FOSP melanoma cohort.}
